% Generated draft: requires mathematical and source review.
\documentclass[11pt]{article}
\usepackage[margin=25mm]{geometry}
\usepackage{fontspec}
\setmainfont{Latin Modern Roman}
\usepackage{amsmath,amssymb,mathtools,unicode-math}
\setmathfont{Latin Modern Math}
\usepackage{enumitem,needspace,xurl,hyperref}
\setlength{\parindent}{0pt}
\setlength{\parskip}{5pt}
\setlength{\emergencystretch}{3em}
\newcommand{\N}{\mathbb{N}}
\newcommand{\Z}{\mathbb{Z}}
\newcommand{\Q}{\mathbb{Q}}
\newcommand{\R}{\mathbb{R}}
\newcommand{\C}{\mathbb{C}}
\newcommand{\F}{\mathbb{F}}
\newcommand{\A}{\mathbb{A}}
\newcommand{\PP}{\mathbb{P}}
\newcommand{\E}{\mathbb{E}}
\newcommand{\eps}{\varepsilon}
\newcommand{\dd}{\,\mathrm{d}}
\begin{document}
\section*{UnsolvedMath ID: 30005382}
\textbf{Source problem code:} catalogue entry
\subsection{Definitions and mathematical statement}
The supplied question concerns an object denoted by \(J\cdot K\), described in the registry as a knot Floer module associated with a knot \(K\). The record does not define \(J\), the operation \(\cdot\), the coefficient ring, or the module structure. Accordingly, write \(M(K)=J\cdot K\) only as a placeholder for the intended construction.

The following is a conditional formulation of the requested construction. Let \(\mathcal C\) denote the intended knot concordance group, with operation induced by connected sum \(\#\), and let \(A\) be an abelian group written additively. Seek a rule \(F\), extracted directly from the objects \(M(K)\), for which
\[
\Phi\colon\mathcal C\longrightarrow A,\qquad [K]\longmapsto F(M(K))
\]
is well defined and is a group homomorphism. Explicitly, for all knots \(K,L\) in the intended category, require
\[
[K]=[L]\ \Longrightarrow\ F(M(K))=F(M(L)),
\qquad
F(M(K\# L))=F(M(K))+F(M(L)).
\]
The source asks that the extraction be performed without first tensoring \(J\cdot K\) with a simpler module.

This template is not a fully specified mathematical problem: the ambient knot category, notion of concordance, construction \(M\), target \(A\), permitted dependence of \(F\) on module data, and meaning of a ``simpler module'' are absent. The source also supplies no nontriviality or normalization requirement; the displayed algebraic conditions alone allow the zero homomorphism. No additional requirement is imposed here. The supplied background identifies indecomposable complexes over a non-principal-ideal domain as an obstacle, but does not specify the ring, complexes, or their relation to \(J\cdot K\).
\subsection{Short English statement}
The question asks how to obtain a knot concordance homomorphism directly from the object \(J\cdot K\), without first simplifying it by tensoring with another module. The supplied record leaves the underlying construction, concordance category, target group, and intended nontriviality unspecified, so the mathematical formulation remains conditional.
\subsection{Sources}
\begin{itemize}
\item Supplied local UnsolvedMath registry record, permanent ID \(30005382\), title \emph{Concordance Homomorphisms from Knot Floer Modules}, exact source problem code \texttt{catalogue entry}, in the Oberwolfach Reports Open Problems collection. The proposed year is recorded as 2023; no individual proposer is supplied.
\item Supplied report citation: \emph{Morphisms in Low Dimensions} (2023), Oberwolfach Reports, DOI 10.4171/owr/2023/4, pp. 215--259. Literature-source label: Oberwolfach Report \(4/2023\). \url{https://doi.org/10.4171/owr/2023/4}. No specific contribution, passage, or problem page is supplied.
\item Supplied literature reference: Schuetz, \emph{Integral Rasmussen invariant}. \url{https://arxiv.org/abs/2202.00445}. The record does not explain the precise connection of this reference to the construction \(J\cdot K\).
\item Literature assessment dated exactly \texttt{2026-08-22}: the supplied curation labels the problem open, describes obstacles involving indecomposable complexes over a non-principal-ideal domain, and treats extraction of a direct concordance homomorphism as unresolved. This is supplied dated curation derived from the Oberwolfach report, not an independently verified literature assessment. The source itself requests independent verification of time-sensitive status and cited claims.
\item The registry labels the statement \texttt{corrected\_verified} and reports that it was checked against the official Oberwolfach report. These are supplied verification claims, not independent verification performed for this draft.
\end{itemize}
\end{document}
